\section{Review Method}

The Introduction framed AI-assisted development frameworks as a layer that organizes the work of whoever uses a development agent, and argued that comparing them requires examining the function of each mechanism rather than surface metrics such as command counts or popularity. We investigated this object through a systematic literature review (SLR) conducted under the Kitchenham protocol \cite{kitchenham2009slr}. The review covers only academic literature, that is, peer-reviewed papers and archived preprints with verifiable bibliographic records; sources without a citable paper were out of scope by construction. The cutoff date was May 26, 2026, and the coverage window ran from 2018 to 2026.

The protocol was fixed before execution and defines the research questions, the sources consulted, the inclusion and exclusion criteria, the extraction schema, and the snowballing strategy. The corpus grounds the analytical dimensions of the taxonomy and establishes the state of knowledge on agents, specification, validation, and risk. Figure~\ref{fig:metodologia} summarizes the review flow, from identification to reporting, with the real counts at each step.

\begin{figure}[h!]
\centering
\includegraphics[width=0.92\textwidth]{methodology_en.pdf}
\caption{Systematic literature review flow, from the Kitchenham protocol to reporting, with the real counts at each step.}
\label{fig:metodologia}
\end{figure}

The review is guided by the five research questions presented in the Introduction (RQ1 through RQ5). RQ5, about the gaps that remain for comparing, evaluating, and governing these frameworks in real projects, is a synthesis-level question answered transversally in the Research Agenda, which is why the extraction matrix records only RQ1 through RQ4 per study.

The search was automated across arXiv, the ACM Digital Library, IEEE Xplore, Crossref, OpenAlex, and Semantic Scholar, with Scopus and Web of Science reached through open metadata via OpenAlex and Google Scholar used for citation checking. The descriptors were organized into four conceptual blocks (AI-assisted development, agentic software engineering, specification as artifact, and organization by workflows), with Boolean strings adapted to each database and documented in the protocol. In line with standard systematic-review practice, the candidate set also received one directed seed: a study within the initial scope, known to the authors and bearing a citable paper, that the automatic retrieval had missed.

The eligibility criteria were fixed in the protocol and are summarized in Table~\ref{tab:criterios}. The decision rule is aggregate: a source that fails criterion I1 is excluded; among sources that satisfy I1, only those that satisfy at least two of I2, I3, and I4 are included. Exclusion reasons were recorded operationally for each discarded source, allowing an auditable account of the discards in each pass.

\begin{table}[h!]
\centering
\caption{Inclusion and exclusion criteria applied in study selection.}
\label{tab:criterios}
\small
\begin{tabularx}{\textwidth}{@{}p{0.46\textwidth}Y@{}}
\toprule
\textbf{Inclusion criteria} & \textbf{Exclusion criteria} \\
\midrule
I1: describes a framework, method, toolkit, platform, IDE, CLI, or kit that organizes AI development across more than one step. & E1: addresses only language models, APIs, or code benchmarks, with no process layer. \\
I2: describes production or consumption of persistent artifacts (specs, plans, PRDs, tasks, logs, evidence). & E2: describes a generic code assistant, with no persistent artifacts, roles, commands, or validation of its own. \\
I3: defines roles, agents, personas, skills, workflows, or coordination commands. & E3: discusses LLMs for a single isolated task, with no framework or process proposal. \\
I4: describes integration with code execution, tests, terminal, IDE, browser, or repositories. & E4: is purely promotional, with no criteria, architecture, or traceable evidence. \\
I5: is an academic paper or archived preprint that is technically verifiable and citable. & E5: is a duplicate of a more complete or more recent source. \\
I6: available in English or Portuguese. & E6: has no citable paper, or is inaccessible for verification. \\
I7: published or updated between 2018 and 2026. & E7: is only an abstract, slide, brief announcement, or opinion without technical detail. \\
I8: allows at least four fields of the extraction schema to be extracted. & E8: mentions the frameworks only tangentially, with no substantive analysis. \\
\bottomrule
\end{tabularx}
\end{table}

The multi-database search returned 100 raw records, reduced to 85 candidates after deduplication by DOI and title, the 2018 to 2026 temporal filter, and the one directed seed described above. These 85 candidates went through a first title-and-abstract screening, which advanced 28 studies to full-text reading and excluded 57. In the second pass, the 28 full texts were read in full under the same criteria, yielding 25 included studies and 3 excluded. The counts and exclusion reasons for both passes were recorded in an append-only log for auditing.

From the 25 included studies, we applied backward and forward snowballing following Wohlin's guidelines \cite{wohlin2014snowballing}, with one depth level and a limit of three rounds, each seeded by the previous round's inclusions. The three rounds generated 69, 50, and 30 new candidates, respectively; of these, 9, 5, and 4 advanced to full text and 6, 3, and 3 were included, and the forward search returned no eligible citations in any round. In total, 149 additional candidates were examined and 12 new studies incorporated, consolidating the final corpus at 37 included studies. One source approved at the third-round screening lacked full text at the cutoff date and was kept as an operational pending item, not an exclusion.

Ten structured fields were extracted from each included study: name, nature, origin, artifacts, roles, execution, validation, portability, evidence, and risks. Consolidating the extraction into a matrix allowed the studies to be compared uniformly. The synthesis followed a hybrid thematic approach: inductive themes were identified from the extraction records and then confronted with the six-dimension taxonomy proposed in the review (specification, context, roles, execution, validation, and portability), reporting the overlap, the thinly covered dimensions, and the themes that emerge outside the taxonomy.

Finally, note that a formal quality assessment by checklist was not applied in this round: the strength of the evidence is classified by source category, distinguishing peer-reviewed study, conceptual paper, and archived preprint. Because the field moves quickly, observations about a study's reported versions, integrations, or tooling should be read as a snapshot of the corpus consulted at the cutoff date. The remaining limitations that bound the reading of the results are consolidated in the Review Limitations section.
